
\definecolor{actives}{HTML}{000957}
\definecolor{reactives}{HTML}{9B0000}
\definecolor{pasives}{HTML}{3F0071}
\definecolor{hybrids}{HTML}{FFD523}

\tikzset{
    parent/.style          = {align = center, text width = 2.5cm, rounded corners = 3pt},
    child/.style           = {align = center, text width = 2.5cm, rounded corners = 3pt},
    grandchild/.style      = {align = center, text width = 4.0cm, rounded corners = 3pt},
    greatgrandchild/.style = {align = center, text width = 5.0cm, rounded corners = 3pt},
    referenceblock/.style  = {align = center, text width = 1.5cm, rounded corners = 2pt}
}

\resizebox{\textwidth}{!}{
    \begin{forest}
        for tree = {
            forked edges,
            grow' = 0,
            draw,
            rounded corners,
            node options = {align = center,},
            text width = 2.7cm,
        },
        [Interfaces cerebro computadora, fill = black, text = white, parent
            [Activas, for tree = {fill = actives!45, child}
                [Movimientos Imaginarions, fill = actives!30, grandchild
                    [Potenciales Corticales Lentos, fill = actives!20, greatgrandchild]
                    [Potenciales Corticales Relacionados a Movimientos, fill = actives!20, greatgrandchild]
                    [Nivel de Sincronización Neuronal, fill = actives!20, greatgrandchild]
                    [Potencia de las Señales EEG, fill = actives!20, greatgrandchild]
                ]
                [Rotación Mental de Cuerpos 3D, fill = actives!30, grandchild]
                [Lingüística Mental, fill = actives!30, grandchild]
                [Musicalización Interna, fill=actives!30, grandchild]
            ]
            [Reactivas, for tree = {fill = reactives!45, child}
                [P300, fill = reactives!30, grandchild]
                [Potenciales Evocados, fill = reactives!30, grandchild
                    [Visuales, fill = reactives!20, greatgrandchild]
                    [Auditivos, fill = reactives!20, greatgrandchild]
                    [Somato sensoriales, fill = reactives!20, greatgrandchild]
                ]
            ]
            [Pasivas, for tree = {fill = pasives!45, child}
                [Ansiedad, fill = pasives!30, grandchild]
                [Fatiga, fill = pasives!30, grandchild]
                [Frustración, fill = pasives!30, grandchild]
                [Atención, fill = pasives!30, grandchild]
            ]
            [Híbridas, for tree = {fill = hybrids!45, child}
                [Combinación de Sistemas, fill = hybrids!30, grandchild
                    [Movimientos Imaginarios y P300, fill = hybrids!20, greatgrandchild]
                    [Potenciales Evocados Visuales y Tareas Mentales, fill = hybrids!20, greatgrandchild]
                ]
                [Combinación de Señales, fill = hybrids!30, grandchild
                    [EEG Y ECG, fill = hybrids!20, greatgrandchild]
                    [EEG y EMG, fill = hybrids!20, greatgrandchild]
                ]
            ]
        ]
    \end{forest}
}

% \definecolor{main_idea}{HTML}{3E8E7E}
% \definecolor{second_idea}{HTML}{7CD1B8}
% \definecolor{third_idea}{HTML}{FABB51}
% \definecolor{fourth_idea}{HTML}{FAEDC6}

% \resizebox{0.85\textwidth}{!}{
% \begin{tikzpicture}[
%     mindmap,
%     grow cyclic,
%     every node/.style = {concept, circular drop shadow},
%     level 1 concept/.append style = {sibling angle = 360/4, level distance = 10em},
%     level 2 concept/.append style = {sibling angle = 45},
%     level 3 concept/.append style = {sibling angle = 40},
%     extra concept/.append style = {text = black}]

%     % BCI Systems - Main Idea
%     \path[mindmap, concept color = main_idea, text = black] node[concept] {Interfaces cerebro computadora} [clockwise from = 45]
    
%         % Active systems
%         child[concept color = second_idea] { node[concept] {Activos} [clockwise from = 130]
%             child[concept color = third_idea] { node[concept] {Movimientos Imaginarios} [clockwise from = -120]
%                 child[concept color = fourth_idea] { node[concept] {Potenciales Corticales Lentos} }
%                 child[concept color = fourth_idea] { node[concept] {Potenciales Corticales Relacionados al Movimiento} }
%                 child[concept color = fourth_idea] { node[concept] {Nivel de Sincronización Neuronal} }
%                 child[concept color = fourth_idea] { node[concept] {Potencia de las señales EEG} }
%             }
%             child[concept color = third_idea] { node[concept] {Rotación Mental de Cuerpos 3D} }
%             child[concept color = third_idea] { node[concept] {Lingüística Mental} }
%             child[concept color = third_idea] { node[concept] {Musicalización Interna} }
%         }
%         child[concept color = second_idea] { node[concept] {Reactivos} [clockwise from = 45]
%             child[concept color = third_idea] { node[concept] {P300} }
%             child[concept color = third_idea] { node[concept] {Potenciales Evocados en Estado Estable} }
%             child[concept color = third_idea] { node[concept] {Potenciales Evocados} }
%         }
%         child[concept color = second_idea] { node[concept] {Pasivos} [clockwise from = -20]
%             child[concept color = third_idea] { node[concept] {Ansiedad} }
%             child[concept color = third_idea] { node[concept] {Fatiga} }
%             child[concept color = third_idea] { node[concept] {Frustración} }
%             child[concept color = third_idea] { node[concept] {Atención} }
%         }
%         child[concept color = second_idea] { node[concept] {Híbridos} [clockwise from = -135]
%             child[concept color = third_idea] { node[concept] {Combinación de sistemas} [clockwise from = -45]
%                 child[concept color = fourth_idea] { node[concept] {Movimientos Imaginarios y P3000} }
%                 child[concept color = fourth_idea] { node[concept] {Potenciales Evocados Visuales y Tareas Mentales} }
%             }
%             child[concept color = third_idea] { node[concept] {Combinación de señales} [clockwise from = 90]
%                 child[concept color = fourth_idea] { node[concept] {EEG Y ECG} }
%                 child[concept color = fourth_idea] { node[concept] {EEG y EMG} }
%             }
%         }
%     ;
% \end{tikzpicture}
% }